\section{Keep a complete product before taking its remainder}
For the scalar function, both residues are at most \(2^{32}-1\). Their product is
less than \(2^{64}\), so a 64-bit low product contains all its bits. This is why the
word-sized interface uses \texttt{dmulu} and \texttt{dmodu}. A word multiply would
discard bits needed by the remainder. A signed doubleword remainder would give a
different numerical interpretation when the product's bit 63 is one.

Take base \texttt{fffffffe} and modulus \texttt{ffffffff}, both unsigned 32-bit
values. Squaring the base gives hexadecimal \texttt{fffffffc00000004}, or decimal
18446744056529682436. Its unsigned remainder modulo 4294967295 is one. Interpreting
the same product bits as a signed 64-bit integer instead gives \(-17179869180\),
which equals \(-4\cdot4294967295\); that signed remainder is zero. Merely checking
small positive products would miss the error.

Load a limb with \texttt{lwu} when it is to be an unsigned magnitude in a 64-bit
calculation. A load with \texttt{lw} sign-extends bit 31. Both loads produce the same
low word, but the full-register multiplication receives different operands. Likewise,
\texttt{addu} produces a word result: it cannot retain a carry into bit 32. Use
\texttt{daddu} for the widened sums below.

Do not extend this scalar strategy by simply replacing its inputs with 128-bit
objects and keeping the low half of their product. For example,
\(2^{127}\cdot2\bmod(2^{128}-1)=1\). Truncating the product to 128 bits first gives
zero, whose remainder is zero. A later division cannot recover discarded bits.
